% GENERATED FILE. Edit canonical lessons, then run npm run generate:books.
% canonical-source-hash: fnv1a64:c43cc9b79151dfd6
% canonical-lessons: ML-C447-netuka, ML-C447-gunanilavaram, ML-C447-parajayappetuka, ML-C447-laksyamituka, ML-C447-karyaksamata

\chapter{Aims and Standards}
\label{ch:ml-aims-and-standards}
\hlchaptermodality{447}

\begin{chapteropening}
\textbf{By the end of this chapter:} \emph{I can say to achieve, a quality standard, to fail, to aim at and efficiency.}

Complete the last lesson of chapter 447: I can say to achieve, a quality standard, to fail, to aim at and efficiency.
\end{chapteropening}

\section[\texorpdfstring{\ml{നേടുക} (nēṭuka)}{നേടുക (nēṭuka)}]{\ml{നേടുക} (nēṭuka) — to achieve; to gain}

\label{lesson:ML-C447-netuka}



\begin{quote}
Before the new one: say the Malayalam for logical, then the Malayalam for reliable.
\end{quote}

\subsection*{You'll want to know: \ml{നേടുക}}
\textbf{\ml{നേടുക}} — \emph{nēṭuka} — \textquotedblleft{}to achieve; to gain\textquotedblright{}.

\begin{quote}
\ml{അവൾ} \ml{പരീക്ഷയിൽ} \ml{ഒന്നാം} \ml{സ്ഥാനം} \ml{നേടി}.
\end{quote}

\begin{quote}
\emph{avaḷ parīkṣayil onnāṁ sthānaṁ nēṭi.}
\end{quote}

\begin{quote}
She achieved first place in the exam.
\end{quote}

Past: \ml{നേടി}, \emph{nēṭi}.

Its past is \ml{നേടി}: \ml{ഒന്നാം} \ml{സ്ഥാനം} \ml{നേടി}, won first place.

\begin{cousinweb}
\textbf{\ml{നേടുക}} (\emph{nēṭuka}) is native Dravidian: DEDR 3766 lists Malayalam \emph{nēṭuka}, \textquotedblleft{}to obtain, to get\textquotedblright{}, with Tamil \emph{nēṭu}, \textquotedblleft{}to seek, to earn\textquotedblright{}.
\end{cousinweb}

\begin{grammarlens}[title={What you've built}]
One more word for saying what happens if.
\end{grammarlens}

\subsection*{Your turn}
\begin{itemize}
\raggedright
  \item \emph{Say it:} \emph{nēṭuka}
  \item \emph{Say it:} \emph{nēṭuka}, once more
  \item \emph{Say it:} \emph{avaḷ parīkṣayil onnāṁ sthānaṁ nēṭi.} — She achieved first place in the exam.
  \item \emph{Recall:} say the Malayalam for logical, then the Malayalam for reliable, then say \emph{nēṭuka} again
\end{itemize}

\begin{tcolorbox}[breakable,colback=teal!4,colframe=teal!35!black,title={Before you move on}]
What does \ml{നേടുക} mean? (\textquotedblleft{}To achieve; to gain\textquotedblright{}.) Say it once more.
\end{tcolorbox}
\section[\texorpdfstring{\ml{ഗുണനിലവാരം} (guṇanilavāraṁ)}{ഗുണനിലവാരം (guṇanilavāraṁ)}]{\ml{ഗുണനിലവാരം} (guṇanilavāraṁ) — quality standard; quality}

\label{lesson:ML-C447-gunanilavaram}



\begin{quote}
Before the new one: say the Malayalam for reliable, then the Malayalam for to achieve.
\end{quote}

\subsection*{You'll want to know: \ml{ഗുണനിലവാരം}}
\textbf{\ml{ഗുണനിലവാരം}} — \emph{guṇanilavāraṁ} — \textquotedblleft{}quality standard; quality\textquotedblright{}.

\begin{quote}
\ml{ആ} \ml{കടയിലെ} \ml{സാധനങ്ങൾക്ക്} \ml{നല്ല} \ml{ഗുണനിലവാരം} \ml{ഉണ്ട്}.
\end{quote}

\begin{quote}
\emph{ā kaṭayile sādhanaṅṅaḷkkŭ nalla guṇanilavāraṁ uṇṭŭ.}
\end{quote}

\begin{quote}
The goods in that shop are of a good standard.
\end{quote}

\ml{ഗുണം}, quality, and \ml{നിലവാരം}, a level: the level of quality.

\begin{cousinweb}
\textbf{\ml{ഗുണനിലവാരം}} (\emph{guṇanilavāraṁ}) is a Malayalam compound: \ml{ഗുണം}, \textquotedblleft{}quality\textquotedblright{}, from Sanskrit \emph{guṇa} (CDIAL 4190), and \ml{നിലവാരം}, \textquotedblleft{}a level, a standard\textquotedblright{}, built on native \ml{നില}, \textquotedblleft{}a standing\textquotedblright{} (DEDR 3675); the -\ml{വാരം} part is not settled.
\end{cousinweb}

\begin{grammarlens}[title={What you've built}]
One more word for saying what happens if.
\end{grammarlens}

\subsection*{Your turn}
\begin{itemize}
\raggedright
  \item \emph{Say it:} \emph{guṇanilavāraṁ}
  \item \emph{Say it:} \emph{guṇanilavāraṁ}, once more
  \item \emph{Say it:} \emph{ā kaṭayile sādhanaṅṅaḷkkŭ nalla guṇanilavāraṁ uṇṭŭ.} — The goods in that shop are of a good standard.
  \item \emph{Recall:} say the Malayalam for reliable, then the Malayalam for to achieve, then say \emph{guṇanilavāraṁ} again
\end{itemize}

\begin{tcolorbox}[breakable,colback=teal!4,colframe=teal!35!black,title={Before you move on}]
What does \ml{ഗുണനിലവാരം} mean? (\textquotedblleft{}Quality standard; quality\textquotedblright{}.) Say it once more.
\end{tcolorbox}
\section[\texorpdfstring{\ml{പരാജയപ്പെടുക} (parājayappeṭuka)}{പരാജയപ്പെടുക (parājayappeṭuka)}]{\ml{പരാജയപ്പെടുക} (parājayappeṭuka) — to be defeated; to fail}

\label{lesson:ML-C447-parajayappetuka}



\begin{quote}
Before the new one: say the Malayalam for to achieve, then the Malayalam for quality standard.
\end{quote}

\subsection*{You'll want to know: \ml{പരാജയപ്പെടുക}}
\textbf{\ml{പരാജയപ്പെടുക}} — \emph{parājayappeṭuka} — \textquotedblleft{}to be defeated; to fail\textquotedblright{}.

\begin{quote}
\ml{ആദ്യ} \ml{ശ്രമത്തിൽ} \ml{അവൻ} \ml{പരാജയപ്പെട്ടു}.
\end{quote}

\begin{quote}
\emph{ādya śramattil avan parājayappeṭṭu.}
\end{quote}

\begin{quote}
He failed at the first attempt.
\end{quote}

Past: \ml{പരാജയപ്പെട്ടു}, \emph{parājayappeṭṭu}.

The noun is \ml{പരാജയം}, a defeat, a failure.

\begin{cousinweb}
\textbf{\ml{പരാജയപ്പെടുക}} (\emph{parājayappeṭuka}) is built from Sanskrit \emph{parājaya}, \textquotedblleft{}defeat, loss\textquotedblright{} (Monier-Williams), with -\ml{പ്പെടുക} from native \ml{പെടുക} (DEDR 3853): to fall into defeat.
\end{cousinweb}

\begin{grammarlens}[title={What you've built}]
One more word for saying what happens if.
\end{grammarlens}

\subsection*{Your turn}
\begin{itemize}
\raggedright
  \item \emph{Say it:} \emph{parājayappeṭuka}
  \item \emph{Say it:} \emph{parājayappeṭuka}, once more
  \item \emph{Say it:} \emph{ādya śramattil avan parājayappeṭṭu.} — He failed at the first attempt.
  \item \emph{Recall:} say the Malayalam for to achieve, then the Malayalam for quality standard, then say \emph{parājayappeṭuka} again
\end{itemize}

\begin{tcolorbox}[breakable,colback=teal!4,colframe=teal!35!black,title={Before you move on}]
What does \ml{പരാജയപ്പെടുക} mean? (\textquotedblleft{}To be defeated; to fail\textquotedblright{}.) Say it once more.
\end{tcolorbox}
\section[\texorpdfstring{\ml{ലക്ഷ്യമിടുക} (lakṣyamiṭuka)}{ലക്ഷ്യമിടുക (lakṣyamiṭuka)}]{\ml{ലക്ഷ്യമിടുക} (lakṣyamiṭuka) — to aim at; to target}

\label{lesson:ML-C447-laksyamituka}



\begin{quote}
Before the new one: say the Malayalam for quality standard, then the Malayalam for to be defeated.
\end{quote}

\subsection*{You'll want to know: \ml{ലക്ഷ്യമിടുക}}
\textbf{\ml{ലക്ഷ്യമിടുക}} — \emph{lakṣyamiṭuka} — \textquotedblleft{}to aim at; to target\textquotedblright{}.

\begin{quote}
\ml{കമ്പനി} \ml{അടുത്ത} \ml{വർഷം} \ml{ഇരട്ടി} \ml{ലാഭം} \ml{ലക്ഷ്യമിടുന്നു}.
\end{quote}

\begin{quote}
\emph{kampani aṭutta varṣaṁ iraṭṭi lābhaṁ lakṣyamiṭunnu.}
\end{quote}

\begin{quote}
The company is aiming at double the profit next year.
\end{quote}

Past: \ml{ലക്ഷ്യമിട്ടു}, \emph{lakṣyamiṭṭu}.

Built on \ml{ലക്ഷ്യം}, an aim, with \ml{ഇടുക}, to put.

\begin{cousinweb}
\textbf{\ml{ലക്ഷ്യമിടുക}} (\emph{lakṣyamiṭuka}) is a Malayalam compound: \ml{ലക്ഷ്യം}, \textquotedblleft{}an aim\textquotedblright{}, from Sanskrit \emph{lakṣya}, \textquotedblleft{}a mark, a target\textquotedblright{} (Monier-Williams), and \ml{ഇടുക}, \textquotedblleft{}to put\textquotedblright{}, native Dravidian (DEDR 442): to set one's mark.
\end{cousinweb}

\begin{grammarlens}[title={What you've built}]
One more word for saying what happens if.
\end{grammarlens}

\subsection*{Your turn}
\begin{itemize}
\raggedright
  \item \emph{Say it:} \emph{lakṣyamiṭuka}
  \item \emph{Say it:} \emph{lakṣyamiṭuka}, once more
  \item \emph{Say it:} \emph{kampani aṭutta varṣaṁ iraṭṭi lābhaṁ lakṣyamiṭunnu.} — The company is aiming at double the profit next year.
  \item \emph{Recall:} say the Malayalam for quality standard, then the Malayalam for to be defeated, then say \emph{lakṣyamiṭuka} again
\end{itemize}

\begin{tcolorbox}[breakable,colback=teal!4,colframe=teal!35!black,title={Before you move on}]
What does \ml{ലക്ഷ്യമിടുക} mean? (\textquotedblleft{}To aim at; to target\textquotedblright{}.) Say it once more.
\end{tcolorbox}
\section[\texorpdfstring{\ml{കാര്യക്ഷമത} (kāryakṣamata)}{കാര്യക്ഷമത (kāryakṣamata)}]{\ml{കാര്യക്ഷമത} (kāryakṣamata) — efficiency}

\label{lesson:ML-C447-karyaksamata}



\begin{quote}
Before the new one: say the Malayalam for to be defeated, then the Malayalam for to aim at.
\end{quote}

\subsection*{You'll want to know: \ml{കാര്യക്ഷമത}}
\textbf{\ml{കാര്യക്ഷമത}} — \emph{kāryakṣamata} — \textquotedblleft{}efficiency\textquotedblright{}.

\begin{quote}
\ml{പുതിയ} \ml{മാനേജർ} \ml{ഓഫീസിന്റെ} \ml{കാര്യക്ഷമത} \ml{കൂട്ടി}.
\end{quote}

\begin{quote}
\emph{putiya mānējar ōphīsinṟe kāryakṣamata kūṭṭi.}
\end{quote}

\begin{quote}
The new manager raised the office's efficiency.
\end{quote}

The adjective is \ml{കാര്യക്ഷമമായ}, efficient.

\begin{cousinweb}
\textbf{\ml{കാര്യക്ഷമത}} (\emph{kāryakṣamata}) is a learned compound built from Sanskrit \emph{kārya}, \textquotedblleft{}business, work\textquotedblright{} (CDIAL 3078), and \emph{kṣamatā}, \textquotedblleft{}ability, fitness\textquotedblright{} (Monier-Williams): fitness for the work.
\end{cousinweb}

\begin{grammarlens}[title={What you've built}]
One more word for saying what happens if.
\end{grammarlens}

\subsection*{Your turn}
\begin{itemize}
\raggedright
  \item \emph{Say it:} \emph{kāryakṣamata}
  \item \emph{Say it:} \emph{kāryakṣamata}, once more
  \item \emph{Say it:} \emph{putiya mānējar ōphīsinṟe kāryakṣamata kūṭṭi.} — The new manager raised the office's efficiency.
  \item \emph{Recall:} say the Malayalam for to be defeated, then the Malayalam for to aim at, then say \emph{kāryakṣamata} again
\end{itemize}

\begin{tcolorbox}[breakable,colback=teal!4,colframe=teal!35!black,title={Before you move on}]
What does \ml{കാര്യക്ഷമത} mean? (\textquotedblleft{}Efficiency\textquotedblright{}.) Say it once more.
\end{tcolorbox}
